% ECONOMICS_PROBLEM: {"id": "C21-453", "title": "External validity and transportability", "jel_code": "C21", "source_id": "OP-0102"}
% !TeX program = xelatex
\documentclass[11pt,a4paper]{article}
\usepackage[margin=23mm]{geometry}
\usepackage{fontspec}
\setmainfont{Latin Modern Roman}
\usepackage{amsmath,amssymb,mathtools}
\usepackage{xcolor,enumitem,booktabs,longtable,array,xurl,hyperref}
\definecolor{muted}{HTML}{53636C}
\hypersetup{colorlinks=true,urlcolor=blue,linkcolor=blue}
\setlength{\parindent}{0pt}
\setlength{\parskip}{5pt}
\newcommand{\R}{\mathbb R}
\newcommand{\E}{\mathbb E}
\newcommand{\N}{\mathbb N}
\newcommand{\ind}{\mathbf 1}
\newcommand{\Var}{\operatorname{Var}}
\newcommand{\Cov}{\operatorname{Cov}}
\newcommand{\supp}{\operatorname{supp}}
\DeclareMathOperator*{\argmin}{arg\,min}
\DeclareMathOperator*{\argmax}{arg\,max}
\newcommand{\entrymeta}[3]{{\sffamily\small\color{muted}\textbf{\texttt{#1}}\quad|\quad #2\par #3\par}}
\newcommand{\Problemref}[1]{\texttt{#1}}
\newcommand{\JELref}[2]{\texttt{#2} (JEL \texttt{#1})}
\begin{document}
\section*{External validity and transportability}
\textbf{Problem ID:} \texttt{C21-453}\quad \textbf{JEL:} \texttt{C21}\par
% BEGIN ECONOMICS STATEMENT
\label{problem:OP-0102}
\label{atlas:prob:E2}

\noindent\textbf{Original identifier:} E2 (atlas item 102). \textbf{Field:} Econometrics and causal inference. \textbf{Classification:} editorial research family with established restricted solutions; current open status not certified.

\subsection*{Definitions and assumptions}

Let $S\in\{0,1\}$ index a target population ($S=0$) and a source experiment ($S=1$). The observed experiment consists of independent source draws $(X,D,Y)$ and independent target draws $X$, with sample sizes $n_1,n_0$ tending to infinity at a stated ratio. Treatment $D\in\{0,1\}$ has potential outcomes $Y(0),Y(1)$; consistency gives $Y=Y(D)$. Assume source randomization or $(Y(0),Y(1))\perp D\mid X,S=1$, with $\eta\le\Pr(D=1\mid X,S=1)\le1-\eta$ for fixed $\eta>0$. The target estimand is $\theta_0=\mathbb E[Y(1)-Y(0)\mid S=0]$. Define source conditional effects $\tau_1(x)=\mathbb E[Y\mid D=1,X=x,S=1]-\mathbb E[Y\mid D=0,X=x,S=1]$. Let $F_s$ denote the covariate law conditional on $S=s$, and define $\tau_0(x)=\mathbb E[Y(1)-Y(0)\mid X=x,S=0]$.

\subsection*{Formal research question}

For explicitly specified departures from effect transportability, characterize the sharp set of $\theta_0$ and construct honest confidence sets. Under $\mathbb E[Y(1)-Y(0)\mid X=x,S=0]=\tau_1(x)$ and absolute continuity of the target covariate law $F_0$ with respect to $F_1$, identification gives
\[
\theta_0=\int\tau_1(x)\,dF_0(x).
\]
Let a sensitivity class instead impose $|\tau_0(x)-\tau_1(x)|\le\delta(x)$ on a declared measurable overlap region $\mathcal O$ where the restriction of $F_0$ is absolutely continuous with respect to $F_1$, with supplied $\delta(x)\ge0$, and $Y(d)\in[a,b]$ everywhere, with known finite $a<b$. Determine sharp bounds including $F_0(\mathcal O^c)$, and procedures uniform over sequences with weakening overlap and estimated density ratios. In particular, characterize when honest intervals can shrink and when unobserved target effect modification creates an irreducible identification width.

\subsection*{Interpretation and scope}

Covariate reweighting addresses observed composition only. A treatment bearing the same name can differ across sites in dosage, implementation or equilibrium consequences; treatment-version invariance must therefore be specified. Prospective target trials can test predictions on sampled environments, while leaving untested environments dependent on assumptions. A marginal-treatment-effect analysis requires its own latent-index selection model and policy support, rather than inheriting the transport formula above. A meaningful output distinguishes uncertainty due to sampling, extrapolation outside overlap, and effect-invariance restrictions.

\subsection*{Source audit and status}

[S1] studies prediction across training-program locations, [S2] develops marginal treatment effects for policy evaluation, and [S3] gives a model-dependent approximation to external-validity bias. None asserts that unrestricted transportability follows from an internally valid experiment. The source question is broad; the overlap-and-sensitivity formulation is an editorial research target. These sources establish motivation and solved restricted cases, without certifying current open status.

\subsection*{References}

\begin{enumerate}[label={[\textnormal{S}\arabic*]},leftmargin=*]

\item V. Joseph Hotz, Guido W. Imbens and Julie H. Mortimer (2005). \emph{Predicting the efficacy of future training programs using past experiences at other locations}. Journal of Econometrics 125(1–2), 241–270. \url{https://www.sciencedirect.com/science/article/pii/S0304407604000806}. \textbf{Locator:} Publisher or institutional abstract and bibliographic record. \textbf{Support and limits:} Cross-site program prediction and heterogeneity; does not establish transportability without assumptions.

\item James J. Heckman and Edward Vytlacil (2005). \emph{Structural Equations, Treatment Effects, and Econometric Policy Evaluation}. Econometrica 73(3), 669–738. \url{https://doi.org/10.1111/j.1468-0262.2005.00594.x}. \textbf{Locator:} Publisher or institutional abstract and bibliographic record. \textbf{Support and limits:} Marginal treatment effects and policy evaluation under structural restrictions; target overlap remains substantive.

\item Isaiah Andrews and Emily Oster (2019). \emph{A simple approximation for evaluating external validity bias}. Economics Letters 178, 58–62. \url{https://www.sciencedirect.com/science/article/pii/S0165176519300655}. \textbf{Locator:} Publisher or institutional abstract and bibliographic record. \textbf{Support and limits:} Local approximation to external validity bias under a selection model; not unrestricted identification.

\end{enumerate}

\subsection*{Common conventions: Source mathematical conventions}

Unless an entry states otherwise, agent and firm sets are finite. State and action spaces are measurable, strategies and outcome maps are measurable, probability laws are specified, and expectations exist for the targets under discussion. Infinite-horizon discount factors lie in $(0,1)$ and discounted objectives are finite. Norms, tolerances and complexity input sizes must be specified for computational comparisons. Constants or primitives that appear locally are inputs, not empirical facts asserted by this catalogue.

\textbf{Identification.} A statistical model $m\in\mathcal M$ implies an observable law $P_m$ and target $\theta(m)$. At law $P$, its identified set is
\[
\Theta(P;\mathcal M)=\{\theta(m):m\in\mathcal M,\ P_m=P\}.
\]
Point identification holds when this set is a singleton, or equivalently when $P_{m_1}=P_{m_2}$ implies $\theta(m_1)=\theta(m_2)$ for all compatible models. Sharp bounds mean bounds attained, or approached when the boundary is not attained, by compatible models. Writing this set is a definition, not a solution: the substantive task is to establish informative restrictions and derive or estimate the set. An unrestricted-confounding model may give only trivial bounds.

\textbf{Inference.} Identification, estimability and valid uncertainty quantification are separate questions. Fix $\alpha\in(0,1)$. A confidence procedure $C_n$ over a declared law class $\mathcal P$ has asymptotic uniform coverage at level $1-\alpha$ when
\[
\liminf_{n\to\infty}\inf_{P\in\mathcal P}\Pr_P\{\theta(P)\in C_n\}\geq1-\alpha.
\]
For set-identified targets the coverage event must be defined separately, for example coverage of the full identified set. If an entry uses identification-indexed models instead of uniquely indexed laws, the same coverage convention is applied over those models. Pointwise limits do not establish uniform validity, and asymptotic validity does not guarantee finite-sample precision. Sample sizes, dependence structures and weak-identification sequences are part of each application.

\textbf{Counterfactuals.} $Y_i(a)$ is a potential outcome under a specified intervention, including its timing, implementation and versions. It is not $Y_i$ conditional on voluntarily choosing $a$. A full intervention vector permits interference; shorthand scalar treatment notation does not silently impose no interference. Assignment, selection, attrition, anticipation and measurement must be specified. A claim about an instrument's local effect is not automatically a population policy effect.

\textbf{Economic equilibrium.} A counterfactual model includes best responses, constraints, feasibility, market clearing, state transitions and expectations consistent with the stated information. When several equilibria exist, either a declared selector is an input or the target is set-valued. Comparisons across policy regimes must use the stated selection convention. Reduced-form relationships are not presumed invariant after an equilibrium-changing intervention.

\textbf{Optimization and welfare.} Optimization is over the stated feasible set. If attainment is not established, $\sup$ or $\inf$ and $\varepsilon$-optimal choices replace an asserted maximizer or minimizer. For example, with value $V=\sup_{a\in\mathcal A}W(a)<\infty$, an $\varepsilon$-optimal set is $\{a\in\mathcal A:W(a)\geq V-\varepsilon\}$ for $\varepsilon>0$. Benefits, costs, risk and health or environmental outcomes can be aggregated only using declared units and conversion weights. Interpersonal utility weights, the treatment of fiscal transfers and foreign profits, discounting, and the behavioral welfare criterion are explicit inputs. A comparative-static derivative additionally requires existence and appropriate differentiability of the selected counterfactual.

\textbf{Sources.} Local labels [S1], [S2], and so forth restart in every question. Locators identify the relevant abstract, model, section or result at the level actually checked; a bibliographic or abstract-level verification is not represented as inspection of an entire proof. Working papers are distinguished from later publications and changing versions. Source summaries are paraphrased. All URL checks and editorial formulations refer to the audit date stated above.
% END ECONOMICS STATEMENT
\end{document}
